\section{Risiken und technische Schulden}
\begin{itemize}
  \item Die REST-basierte Choreographie erzeugt temporale Kopplung zwischen aufeinanderfolgenden Services. 
  Für eine produktive Variante ist asynchrones Messaging, z.\,B. Kafka, zu prüfen.
  \item H2 ist für den lokalen Prototyp geeignet, nicht jedoch für eine produktive Datenhaltung mit Hochverfügbarkeit.
  \item Monitoring ist im PoC fachlich sichtbar, aber noch kein vollständiger Observability-Stack mit zentralem Logging, 
  Metriken und Tracing.
  \item Sicherheitsaspekte wie Authentisierung, Autorisierung und feinere API-Verträge sind für einen Produktivbetrieb weiter auszubauen.
\end{itemize}

\newpage

\section{Quellenbasis}
Die Dokumentation stützt sich auf die Übungsblätter Nr.~4, 5 und 6 sowie die finalen Semesterprojekt-Hinweise der Lehrveranstaltung \cite{uebungsblaetter2026} und die Fallstudie \emph{WirSchiffenDas} \cite{fallstudie2026}. Fachlich orientiert sie sich an der Qualitätssicherung für Microservice-Architekturen nach Schirgi und Brenner \cite{schirgi2021quality} sowie an der Taxonomie von Microservice-Architekturen nach Garriga \cite{garriga2018taxonomy}. Die Struktur der Beschreibung folgt dem arc42-Template \cite{starke2017arc42}.

Ergänzende Projektdokumentation befindet sich in \code{docs/arc42.md}, \code{docs/ddd.md}, \code{docs/requirements.md}, \code{docs/architecture.md}, \code{docs/testing.md} und \code{docs/diagrams/} in unserem Branch auf der \href{https://github.com/yunusemresirin/WirSchiffenDas/tree/main/docs}{GitHub-Repository}.
